% from ega3-1-fr.tex
\oldpage[0\textsubscript{III}]{46}
pour $(\sigma,\tau)\leq(\sigma',\tau')$. Alors on vérifie aussitôt que
\[
\mathrm C^\bullet(A,B;\Gamma^\bullet)
=\bigl(\mathrm C^h(A,B;\Gamma^k)\bigr)_{(h,k)\in\mathbb Z\times\mathbb Z}
\]
est un bicomplexe de groupes abéliens, et
\[
\mathrm L^\bullet(A,B;\Gamma^\bullet)
=\bigl(\mathrm L^h(A,B;\Gamma^k)\bigr)
\]
en est un sous-bicomplexe.
\end{env}

\begin{proposition}[11.8.9]
\label{0.11.8.9-fr}
\emph{L'injection canonique
\[
\mathrm L^\bullet(A,B;\Gamma^\bullet)\longrightarrow
\mathrm C^\bullet(A,B;\Gamma^\bullet)
\]
définit un isomorphisme pour la cohomologie de ces deux bicomplexes.}
\end{proposition}

\begin{proof}
Posons $\mathrm C^{\bullet\bullet}=
\mathrm C^\bullet(A,B;\Gamma^\bullet)$ et
$\mathrm L^{\bullet\bullet}=
\mathrm L^\bullet(A,B;\Gamma^\bullet)$ pour simplifier, et notons que
puisque $\mathrm C^{hk}=\mathrm L^{hk}=0$ pour $h<0$, les secondes suites
spectrales de ces bicomplexes sont régulières
(\hyperref[0.11.3.3-fr]{11.3.3}); l'homomorphisme
$\mathrm L^{\bullet\bullet}\to\mathrm C^{\bullet\bullet}$ fournit donc un
morphisme de suites spectrales
\[
{}''E(\mathrm L^{\bullet\bullet})\longrightarrow
{}''E(\mathrm C^{\bullet\bullet})
\]
qui, pour les termes $E_2$, se réduit à
\[
\label{0.11.8.9.1-fr}
\tag{11.8.9.1}
H^p_{\mathrm{II}}(H^q_{\mathrm I}(\mathrm L^{\bullet\bullet}))
\longrightarrow
H^p_{\mathrm{II}}(H^q_{\mathrm I}(\mathrm C^{\bullet\bullet})).
\]
Mais pour tout $k\in\mathbb Z$, il résulte de
(\hyperref[0.11.8.3-fr]{11.8.3}) que l'homomorphisme
$H^q_{\mathrm I}(\mathrm L^{\bullet,k})\to
H^q_{\mathrm I}(\mathrm C^{\bullet,k})$ est bijectif; la conclusion résulte
donc de (\hyperref[0.11.1.5-fr]{11.1.5}).
\end{proof}

\begin{env}[11.8.10]
\label{0.11.8.10-fr}
De même, avec les notations de (\hyperref[0.11.8.6-fr]{11.8.6}), on a des
homomorphismes canoniques de bicomplexes
\[
\mathrm C^\bullet(A,B;\Gamma^\bullet)\longrightarrow
\mathrm P^\bullet(A,B;\Gamma^\bullet)
\]
(avec des notations évidentes), et le même raisonnement que dans
(\hyperref[0.11.8.9-fr]{11.8.9}), basé cette fois sur
(\hyperref[0.11.8.6-fr]{11.8.6}), montre que cet homomorphisme donne encore
un isomorphisme en cohomologie.
\end{env}

\subsection{Un lemme sur les complexes de type fini}
\label{subsection:0.11.9-fr}

\begin{proposition}[11.9.1]
\label{0.11.9.1-fr}
Soient $\mathcal C$ une catégorie abélienne, $K'$ et $K''$ des parties de
l'ensemble des objets de $\mathcal C$, telles que $K''\subset K'$, et
vérifiant les conditions suivantes :
\begin{enumerate}
\item[{\rm(i)}] \emph{Pour tout objet $A'\in K'$ et tout épimorphisme
$u:A\to A'$ dans $\mathcal C$, il existe un objet $B\in K''$ et un morphisme
$v:B\to A$ tels que $u\circ v$ soit un épimorphisme.}
\item[{\rm(ii)}] \emph{Pour tout couple d'objets $A\in K'$, $B\in K'$ et
tout épimorphisme $u:A\to B$, $\operatorname{Ker}(u)$ appartient à $K'$.}
\item[{\rm(iii)}] \emph{Le produit de deux objets de $K''$ appartient à
$K''$.}
\end{enumerate}
\emph{Soit $P_\bullet=(P_i)_{i\in\mathbb Z}$ un complexe dans $\mathcal C$,
tel que $H_i(P_\bullet)\in K'$ pour tout $i$, et qu'il existe $d$ tel que
$H_i(P_\bullet)=0$ pour $i<d$. Alors il existe un complexe
$Q_\bullet=(Q_i)_{i\in\mathbb Z}$ dans $\mathcal C$, tel que $Q_i\in K''$
pour tout $i$ et $Q_i=0$ pour $i<d$, et un morphisme
$u:Q_\bullet\to P_\bullet$ de complexes tel que le morphisme correspondant
$H_\bullet(Q_\bullet)\to H_\bullet(P_\bullet)$ soit un isomorphisme.}
\end{proposition}

\begin{proof}
Démontrons d'abord la conséquence suivante de la propriété (i) :

\emph{(i bis) Soient $u:C\to B$ un épimorphisme dans $\mathcal C$, $A$ un
objet de $K'$, $v:A\to B$ un morphisme dans $\mathcal C$; il existe alors un
objet $D\in K''$, un épimorphisme $u':D\to A$ et un morphisme $v':D\to C$
tels que le diagramme}
\[
\xymatrix{
D\ar[r]^{u'}\ar[d]_{v'} & A\ar[d]^{v}\\
C\ar[r]_{u} & B
}
\tag{11.9.1.1}\label{0.11.9.1.1-fr}
\]
\emph{soit commutatif.}

Considérons en effet le produit fibré $C\times_B A$ dans $\mathcal C$ et les
projections canoniques $p:C\times_B A\to C$, $q:C\times_B A\to A$, rendant
commutatif le diagramme
